
We denote discrete physical time steps with superscripts and pseudo-time steps with subscripts: $\state^{n}$ is a quantity at physical time step $n$ and $\bx_{\tau}$ at pseudo-time $\tau$. We write $X\sim \dist{x}$ and $X\sim \conddist{x}{y}$ when $X$ is distributed according to $P(X)$ and $P(X|Y)$, identifying distributions through their densities.

\subsection{Bayesian inverse problems and data assimilation}
We outline the data assimilation setting; for more detail see \cite{carrassi_data_2018, evensen_data_2022}. We consider the discrete dynamical system:
\begin{subequations} \label{eq:data_assimilation_problem}
\begin{align}
    \state^{n} &= \forwardmodel(\state^{n-1}) + \modelnoise^{n},  &\modelnoise^{n}\sim \dist[\modelnoise], \label{eq:forward}\\
    \obs^{n} &= \obsoperator(\state^{n}) + \obsnoise^{n},  &\obsnoise^{n}\sim \dist[\obsnoise], \label{eq:observation}
\end{align}
\end{subequations}
where, for $n=1,\ldots,N$, $\state^{n}\in\R^{\statedim}$ is the state, $\obs^n\in\R^{\obsdim}$ the observations, $\modelnoise^n\in\R^{\statedim}$ the model noise, $\obsnoise^n\in\R^{\obsdim}$ the observation noise, $\forwardmodel:\R^{\statedim}\rightarrow\R^{\statedim}$ the forward operator, and $\obsoperator:\R^{\statedim}\rightarrow\R^{\obsdim}$ the observation operator. Eq.~\eqref{eq:forward} induces a prior transition density $\conddist{\state^{n}}{\state^{n-1}}$, and the observation likelihood follows from the noise distribution, $\conddist{\obs^n}{\state^{n}} = \dist[\eta]{\obs^n-\obsoperator(\state^{n})}$.

We are interested in the posterior over trajectories, which by the Markovian property of Eq.~\eqref{eq:data_assimilation_problem} factors autoregressively into the one-step posterior:
\begin{align} \label{eq:posterior_distribution}
    \conddist{\state^{n}}{\obs^n, \state^{n-1}} = \frac{\conddist{\obs^{n}}{\state^{n}}\conddist{\state^{n}}{\state^{n-1}}}{\conddist{\obs^{n}}{\state^{n-1}}},
\end{align}
with $\conddist{\state^{n}}{\state^{n-1}}$ the prior, $\conddist{\obs^{n}}{\state^{n}}$ the likelihood, and the denominator the evidence. We sample the prior with generative models.

\subsection{Interpolation-based generative modelling}
\label{subsec:flow_models}

We build the prior $\conddist{\state^{n}}{\state^{n-1}}$ as a generative model that samples a target conditional $\conddist{\bx_1}{\bx_0}$; in forecasting, $\bx_0=\state^{n-1}$ and $\bx_1=\state^{n}$ recover the prior in Eq.~\eqref{eq:posterior_distribution}. The building block is a process that interpolates, over a pseudo-time $\tau\in[0,1]$ unrelated to physical time, between a tractable source and the target,
\begin{align}\label{eq:general_interpolant}
    \bx_\tau = \underbrace{\alpha_\tau \bx_0 + \sigpath_\tau \latentnoise}_{\text{source}} + \beta_\tau \bx_1,
    \qquad \latentnoise \sim \mathcal{N}(0,\bI),
\end{align}
with $\bx_0$ a deterministic anchor available at sampling time, $\latentnoise$ independent of $\bx_1$, and schedules satisfying $\beta_0=0$, $\beta_1=1$, $\alpha_1=\sigpath_1=0$, so that $\bx_{\tau=1}=\bx_1$. Conditioned on $\bx_1$, Eq.~\eqref{eq:general_interpolant} is Gaussian with mean $\alpha_\tau\bx_0+\beta_\tau\bx_1$ and covariance $\sigpath_\tau^2\bI$, and the marginal $p_\tau(\bx\mid\bx_0)$ is the corresponding Gaussian mixture. Flow matching, stochastic interpolants, and diffusion models are instances of this one path, distinguished only by the schedules and the sampling dynamics.

The path carries two central quantities. The \emph{velocity} transports the source to $\conddist{\bx_1}{\bx_0}$ and is the conditional-mean instantaneous displacement,
\begin{align}\label{eq:general_velocity}
    \fmvel(\bx,\bx_0,\tau)
    = \E\!\left[\dot\alpha_\tau \bx_0+ \dot\sigpath_\tau\latentnoise + \dot\beta_\tau\bx_1 \;\middle|\; \bX_\tau=\bx\right],
\end{align}
so samples follow from the probability-flow ODE $\rd\bX_\tau = \fmvel_\tau(\bX_\tau,\bx_0)\,\rd\tau$ \cite{lipman_flow_2023, albergo_stochastic_2023}. The \emph{score} $\genscore_\tau(\bx,\bx_0) = \nabla_{\bx}\log p_\tau(\bx\mid\bx_0)$ determines the velocity and vice versa (Appendix~\ref{appendix:score_velocity}),
\begin{align}\label{eq:general_score_velocity}
    \genscore_\tau(\bx)
    = \frac{\beta_\tau\,\fmvel(\bx,\bx_0,\tau) - \dot\beta_\tau\,\bx - (\dot\alpha_\tau\beta_\tau - \dot\beta_\tau\alpha_\tau)\,\bx_0}
           {\beta_\tau\,\vscoef_\tau},
    \qquad
    \vscoef_\tau := \frac{\sigpath_\tau(\dot\beta_\tau\sigpath_\tau - \dot\sigpath_\tau\beta_\tau)}{\beta_\tau},
\end{align}
where $\vscoef_\tau$ is the \emph{velocity--score coefficient}. Whichever of the two the network was trained to represent, the other is available in closed form -- the fact that lets one construction below serve all model classes.

\paragraph{Flow matching.}
A pure Gaussian source: $\alpha_\tau\equiv 0$ with $\sigpath_\tau$ the source schedule (e.g.\ $\sigpath_\tau=1-\tau$, $\beta_\tau=\tau$ for rectified flow), and $\state^{n-1}$ entering only as a conditioning input to the network \cite{lipman_flow_2023, liu_rectified_2022}. Training regresses the velocity Eq.~\eqref{eq:general_velocity} without ODE solves; sampling integrates the probability-flow ODE from $\bX_0\sim\mathcal{N}(0,\bI)$.

\paragraph{Stochastic interpolants.}
The anchor is the previous state, $\bx_0=\state^{n-1}$, and noise enters as a Brownian motion $\gamma_\tau\bW_\tau$, $\bW_\tau=\sqrt\tau\,\bz$, $\bz\sim\mathcal{N}(0,\bI)$, so Eq.~\eqref{eq:general_interpolant} reads
\begin{align}\label{eq:stochastic_interpolant}
    \bx_\tau = \alpha_\tau\bx_0 + \beta_\tau\bx_1 + \gamma_\tau\bW_\tau,
\end{align}
with marginal noise scale $\sigpath_\tau=\gamma_\tau\sqrt\tau$ and $\alpha_0=1$, $\gamma_1=0$ \cite{albergo_stochastic_2023, chen_probabilistic_2024}. The Brownian construction admits regressing the drift of the generative SDE directly,
\begin{align}\label{eq:SI_SDE}
    \rd\bX_\tau = \drift_\theta(\bX_\tau,\bx_0,\tau)\,\rd\tau + \gamma_\tau\,\rd\bW_\tau, \quad \bX_0=\bx_0, \quad
    \drift_\theta = \E[\dot\alpha_\tau\bx_0 + \dot\beta_\tau\bx_1 + \dot\gamma_\tau\bW_\tau\mid\bX_\tau=\bx],
\end{align}
again without SDE solves during training; the drift decomposes as $\drift_\theta = \fmvel_\tau + \tfrac12\gamma_\tau^2\genscore_\tau$.

\paragraph{Diffusion models.}
Sampling need not be deterministic: the probability-flow ODE is one member of a family of SDEs with identical marginals \cite{song_score-based_2021, ma_sit_2024}. For any schedule $\gdiff_\tau\ge 0$,
\begin{align}\label{eq:general_sde_family}
    \rd\bX_\tau
    = \big[\fmvel_\theta(\bX_\tau,\bx_0,\tau) + \tfrac{1}{2}\gdiff_\tau^2\,\genscore_\tau(\bX_\tau,\bx_0)\big]\,\rd\tau
    + \gdiff_\tau\,\rd\bW_\tau,
\end{align}
with the score in closed form from Eq.~\eqref{eq:general_score_velocity} (instantiated for both trained objects in Eqs.~\eqref{eq:fm_score} and \eqref{eq:SI_score}, Appendix~\ref{appendix:score_velocity}). Setting $\gdiff_\tau=0$ recovers the flow-matching ODE and $\gdiff_\tau=\gamma_\tau$ the stochastic-interpolant SDE \eqref{eq:SI_SDE}; any $\gdiff_\tau>0$ on an FM prior is exactly a denoising-diffusion (score-SDE) sampler \cite{song_score-based_2021, ho_denoising_2020}. Eq.~\eqref{eq:general_sde_family} is the common interface on which Section~\ref{sec:methodology} operates.
